For the main comparison, we use a fixed-replay benchmark so that differences
primarily reflect the allocation strategy. Each method receives the same $76$
temporally held-out historical routes, case arrivals, and simulator settings.
We repeat the benchmark with seeds $1$--$5$, yielding $380$ fixed-route runs
per method.

The compared methods are R-RRA, R-SHQ, ParkSong-Composite,
Kunkler-Rinderle-Ma, and Batch. They use the same development-selected RF
branching configuration. The processing-time estimates used by ParkSong are
derived exclusively from the training partition.

The BPMN model provides BPMN-valid transition candidates and is also used for
conformance diagnostics. If a diagnostic transition cannot be fired, fixed
replay nevertheless continues along the historical route. Completion therefore
denotes completion of the fixed historical route, not necessarily attainment
of the BPMN final marking.

The experiment uses a $24$-hour arrival window followed by a $60$-day drain
horizon. Cycle time is calculated only for completed routes. Waiting time is
measured from entry into the resource queue until resource assignment, while
throughput and occupation use the full arrival-plus-drain horizon.

Incomplete routes are reported separately; cycle-time results must therefore
be interpreted jointly with completion. Per-seed results and variability are
reported in App.~\ref{app:joao-fixed-replay-details}. The findings are
interpreted descriptively over $76$ routes and five seeded repetitions.